\documentclass[11pt,a4paper]{article}

\usepackage[utf8]{inputenc}
\usepackage[T1]{fontenc}
\usepackage{lmodern}
\usepackage[french]{babel}
\usepackage[margin=2.1cm]{geometry}
\usepackage{microtype}
\usepackage{amsmath,amssymb}
\usepackage{eurosym}
\usepackage{booktabs}
\usepackage{array}
\usepackage[table]{xcolor}
\usepackage{graphicx}
\usepackage{caption}
\usepackage{enumitem}
\usepackage[most]{tcolorbox}
\usepackage{tikz}
\usetikzlibrary{arrows.meta,positioning,calc,fit,backgrounds,decorations.pathreplacing,shapes.geometric}
\usepackage[hidelinks]{hyperref}

\definecolor{navy}{HTML}{1F3864}
\definecolor{bluemid}{HTML}{2E5496}
\definecolor{lightblue}{HTML}{D6E0F0}
\definecolor{accent}{HTML}{C0491F}
\definecolor{vert}{HTML}{3D8361}

\captionsetup{font=small,labelfont={bf,color=navy},labelsep=period}
\graphicspath{{../figures/}}

\newtcolorbox{phrasebox}{colback=vert!7,colframe=vert!70!black,boxrule=0.7pt,
  arc=2pt,left=7pt,right=7pt,top=5pt,bottom=5pt,before skip=6pt,after skip=8pt}
\newtcolorbox{cle}[1][]{colback=lightblue!40,colframe=navy,boxrule=0.6pt,
  arc=2pt,left=7pt,right=7pt,top=5pt,bottom=5pt,#1}

\providecommand{\Prob}{\mathbb{P}}
\providecommand{\E}{\mathbb{E}}
\providecommand{\R}{\mathbb{R}}

\newcommand{\lequoi}[1]{\smallskip\noindent\textcolor{navy}{\textbf{Le problème.}}\ #1\par}
\newcommand{\lidee}[1]{\smallskip\noindent\textcolor{navy}{\textbf{L'idée.}}\ #1\par}
\newcommand{\laformule}[1]{\smallskip\noindent\textcolor{navy}{\textbf{La formule.}}\ #1\par}
\newcommand{\cheznous}[1]{\smallskip\noindent\textcolor{navy}{\textbf{Chez nous.}}\ #1\par}
\newcommand{\lepiege}[1]{\smallskip\noindent\textcolor{accent}{\textbf{Le piège.}}\ #1\par}
\newcommand{\laphrase}[1]{\begin{phrasebox}\textbf{À dire à voix haute.} \emph{#1}\end{phrasebox}}

% Questions du jury : la question, puis la reponse a tenir.
\newtcolorbox{question}[1]{colback=accent!5,colframe=accent!75!black,boxrule=0.6pt,
  arc=2pt,left=7pt,right=7pt,top=5pt,bottom=5pt,before skip=7pt,after skip=4pt,
  fonttitle=\bfseries\color{white},title={#1}}
\newcommand{\rep}[1]{\smallskip\noindent\textcolor{navy}{\textbf{La réponse.}}\ #1\par}
\newcommand{\appui}[1]{\smallskip\noindent\small\textcolor{navy}{\textbf{Appui~:}} #1\par\normalsize}

% Styles des schemas
\tikzset{
  brique/.style={draw=navy,fill=lightblue!45,rounded corners=2pt,inner sep=5pt,
                 align=center,font=\footnotesize,minimum height=9mm},
  pilier/.style={draw=navy,fill=white,circle,inner sep=1.5pt,minimum size=8mm,
                 font=\footnotesize\bfseries,text=navy},
  fleche/.style={-{Stealth[length=2mm]},draw=navy,line width=0.5pt},
  flechef/.style={-{Stealth[length=2.6mm]},draw=accent,line width=1.1pt},
  flechep/.style={-{Stealth[length=1.8mm]},draw=navy!45,line width=0.4pt},
  etiq/.style={font=\scriptsize\itshape,text=navy!75},
}

\setcounter{secnumdepth}{2}
\emergencystretch=1.2em   % noms propres longs (Kolmogorov-Smirnov, Pickands-Balkema)

\title{\vspace{-1.3cm}\color{navy}\bfseries Le mémoire, concept par concept\\[3pt]
  \large Ce que chaque outil résout, comment il marche,
  et la phrase qui l'explique en trente secondes}
\author{Kélian Kaddouri}
\date{21 septembre 2026}

\begin{document}
\maketitle
\vspace{-0.8cm}

\begin{cle}
\textbf{À quoi sert ce document.} Le mémoire empile une vingtaine d'outils, chacun choisi
pour une raison précise, et il est facile de les manipuler correctement sans savoir les
\emph{défendre}. Chaque section suit le même gabarit~: le problème que l'outil résout,
l'idée en langage courant, la formule, ce que ça donne dans notre cas, la phrase à dire à
voix haute quand on vous interroge, et le piège à éviter.

\textbf{La phrase encadrée en vert est le livrable de chaque section.} Si vous ne retenez
qu'elle, vous tenez l'oral.
\end{cle}

\smallskip
\noindent\textcolor{accent}{\textbf{Mise à jour du 21 septembre 2026, à neuf jours du dépôt.}}
Cette révision ajoute ce qui manquait pour tenir l'oral~: une \textbf{partie~X entière de
questions de jury} avec la réponse à tenir et la page d'appui, \textbf{six schémas} qui montrent
les mécanismes au lieu de les décrire, et la mise à jour des nombres. \textcolor{accent}{Deux
nombres de l'annexe étaient périmés et sont corrigés}~: le seuil des excès vaut $20{,}03$~M\euro{}
et non $4{,}176$ (ce dernier est un quantile de table descriptive, jamais un seuil), et le seuil
de robustesse du traité vaut $\kappa^\star=78\,\%$ et non $76$. Le chiffre de tête
($6\,049\to20\,188$~M\euro), les bornes d'identification et le seuil de bascule du classement
entrent également dans l'annexe.

\bigskip
\noindent\textbf{Le schéma d'ensemble.} Cinq étages, de la donnée au capital. Chaque partie de ce
document en traite un.

\begin{center}
\begin{tikzpicture}[node distance=4.5mm]
\node[brique] (d) {\textbf{Données}\\OpRisk, PRC\\Hackmageddon};
\node[brique,right=of d] (f) {\textbf{Fréquence}\\binomiale\\négative};
\node[brique,right=of f] (s) {\textbf{Sévérité}\\corps\,+\,queue GPD\\$\xi=0{,}5954$};
\node[brique,right=of s] (c) {\textbf{Cascade}\\dirigée $W$\\$\rho(W)<1$};
\node[brique,right=of c] (m) {\textbf{Conformité}\\multi-états\\(Markov)};
\node[brique,right=of m,fill=vert!15,draw=vert!70!black] (k)
  {\textbf{Capital}\\$\mathrm{VaR}_{99{,}5\%}(L)$};
\draw[fleche] (d)--(f); \draw[fleche] (f)--(s); \draw[fleche] (s)--(c);
\draw[fleche] (c)--(m); \draw[fleche] (m)--(k);
\draw[flechef] (d.south) to[out=-60,in=-120] node[etiq,below,pos=0.5]
  {la direction de $W$ n'est PAS identifiable ici} (c.south);
\end{tikzpicture}
\end{center}

\noindent La flèche rouge est le sujet du mémoire~: tout le reste se calibre sur la donnée,
\emph{sauf} la direction de la contagion. C'est de là que vient l'identification partielle,
et c'est la question que le jury posera en premier.

\tableofcontents
\newpage

% =====================================================================
\part*{I. D'où viennent les nombres}
\addcontentsline{toc}{part}{I. D'où viennent les nombres}

\section{Les trois sources, et ce que chacune peut dire}

\lequoi{Pour chiffrer un capital, il faut trois choses~: combien d'incidents par an, combien
coûte un incident, et comment les incidents s'enchaînent. Aucune source publique ne porte
les trois.}

\lidee{On ne cherche pas \emph{la} bonne base, on répartit les questions entre les sources
selon ce que chacune peut porter, et on déclare ce qui manque.}

\cheznous{Trois sources, trois rôles disjoints.
\begin{itemize}[itemsep=1pt,topsep=2pt]
\item \textbf{SAS OpRisk Global Data}~: pertes opérationnelles en dollars, catégories de
Bâle, $4\,329$ transitions annuelles consécutives. Elle porte le \textbf{niveau} de
sévérité et le volume. Son pas est annuel, ce qui efface l'ordre des événements.
\item \textbf{Privacy Rights Clearinghouse (PRC)}~: $15\,053$ incidents, mais exprimés en
\emph{nombre d'enregistrements compromis}, pas en euros. Elle sert d'\textbf{axe de
sensibilité} après conversion.
\item \textbf{Sept rapports post-incident}~: peu de volume, mais une enquête y reconstitue
la \emph{séquence} des défaillances. C'est la seule source qui porte l'\textbf{ordre}.
\end{itemize}}

\laphrase{Aucune base ne dit tout. OpRisk donne le niveau et le volume mais pas l'ordre,
PRC donne du volume mais pas d'euros, et les rapports d'enquête donnent l'ordre mais pas le
volume. Je les utilise pour ce que chacune peut porter, et je dis explicitement ce qui
reste absent.}

\lepiege{Ne jamais laisser croire qu'on a \og une base de données DORA\fg{}. Elle n'existe pas.
Le règlement met en place le reporting qui la produira, et c'est justement l'objet de la
recommandation finale du mémoire.}

\section{La conversion Jacobs}

\lequoi{PRC compte des enregistrements compromis, pas des euros. Il faut passer de l'un à
l'autre.}

\lidee{Une régression publiée établit, sur un échantillon d'incidents dont on connaît à la
fois le nombre d'enregistrements et le coût, une relation log-log entre les deux.}

\laformule{$\ln L = 7{,}68 + 0{,}76\,\ln X + \varepsilon$, où $X$ est le nombre
d'enregistrements et $L$ la perte en dollars. Le résidu a un écart-type de $0{,}523$, et
$R^2 = 0{,}512$.}

\cheznous{L'exposant $0{,}76 < 1$ dit que le coût croît \emph{moins vite} que le nombre
d'enregistrements~: un incident dix fois plus gros ne coûte pas dix fois plus cher, mais
environ six fois. Il y a un coût fixe important dans tout incident.}

\laphrase{Jacobs n'est pas un lien fréquence-sévérité, c'est une conversion de sévérité
seule~: elle transforme un nombre d'enregistrements en montant. Et mon capital principal
n'en dépend pas, il est calculé sur OpRisk, qui est déjà en dollars.}

\lepiege{Le modèle traitait cette conversion comme \emph{exacte}, en ignorant le résidu, ce
qui écrasait la dispersion de sévérité. En réinjectant le résidu, le niveau monte de
$7\,\%$ et la queue s'épaissit ($\xi$ passe de $1{,}03$ à $1{,}07$), mais la largeur
relative de l'intervalle ne bouge pas. À dire spontanément~: la conversion était traitée
comme sûre alors qu'elle ne l'est pas.}

% =====================================================================
\section{Le seuil de collecte, et les trois fréquences}

\lequoi{Trois nombres du mémoire portent le nom de \og{}fréquence\fg{} et vont de $0{,}10$ à
$341$ par an. Confondre deux d'entre eux fait dire n'importe quoi.}

\lidee{Ils ne comptent pas la même chose, parce qu'ils ne sont pas lus au même endroit. Une
base ne retient que les pertes au-dessus d'un \textbf{seuil de collecte} : le taux qu'on y lit
est celui des incidents \emph{matériels}, générateurs d'une perte enregistrée, et non celui de
tous les incidents au sens du reporting réglementaire.}

\cheznous{\;\\[-6pt]
\begin{center}\small
\begin{tabular}{@{}l l l@{}}
\toprule
$\lambda_{ref}=341$/an & notifications PRC & \emph{clé de répartition} par vecteur \\
$\lambda=21{,}6$/an & incidents \textsc{tic} matériels, secteur & \textbf{l'entrée du moteur} \\
$\lambda_{\text{entité}}=0{,}10$ à $0{,}21$/an & les mêmes, par firme & l'échelle d'une entité \\
\bottomrule
\end{tabular}
\end{center}}

\laphrase{Le facteur mille six cents entre mes trois fréquences n'est pas une incohérence,
c'est trois niveaux d'observation. Un seul alimente le moteur, les deux autres servent à
répartir et à interpréter. Le seuil de collecte est ce qui les sépare.}

\lepiege{Le taux $\lambda$ est une propriété d'\textbf{entité}, lisible sur un panel
firme-année ; la charge systémique commune, elle, ne se lit qu'au niveau \textbf{agrégé} et
après retrait de tendance. Les confondre revient à imputer au risque systémique ce qui n'est
qu'un effet de taille des firmes.}

\section{Mesurer un biais au lieu de le déclarer}

\lequoi{Tout le monde déclare que sa base est biaisée. Presque personne ne dit de combien. Or
déclarer ne coûte rien, mesurer engage.}

\lidee{Deux biais classiques d'une base de pertes se mesurent avec la base elle-même. Le
\textbf{biais de taille}, parce que la base porte les actifs de la firme sinistrée. La
\textbf{troncature à droite}, parce qu'elle porte la date de saisie en plus de l'année de
survenance.}

\laformule{Pour la taille, une régression log-log $\ln L = a + b\ln(\text{actifs})$ donne
l'élasticité $b$. Mais elle est biaisée par le seuil de collecte : une petite firme n'entre
dans la base que si sa perte est anormalement grande, ce qui \emph{aplatit} la pente. On
corrige par un maximum de vraisemblance sur une lognormale \textbf{tronquée à gauche} en $u_0$,
où la contribution de chaque observation est divisée par $1-\Phi\bigl((\ln u_0-\mu)/s\bigr)$.
Pour la troncature à droite, on construit le triangle survenance $\times$ délai de remontée et
l'on applique un \textbf{chain-ladder} aux comptes.}

\cheznous{Élasticité brute $0{,}042$ avec un $R^2$ de $0{,}005$, ce qui \emph{réfuterait}
l'explication usuelle. Corrigée de la troncature, elle double à $\mathbf{0{,}087}$
$[0{,}026\,;0{,}148]$ et devient significative, mais reste faible : la sévérité transposée à
l'entité notionnelle vaut $0{,}86$ fois la brute, soit $-14\,\%$. Côté remontée, le délai médian
est de \textbf{trois ans} et les $97$ événements observés sur 2023--2025 deviennent
\textbf{$255$} à l'ultime, soit $+163\,\%$.}

\laphrase{Mes deux biais les plus souvent opposés, je ne me contente pas de les signaler : je
les chiffre. Le biais de taille coûte quatorze pour cent, pas un facteur dix. Et la baisse
apparente des incidents après 2022 n'est pas une baisse du risque, c'est une remontée
incomplète que le chain-ladder rattrape.}

\lepiege{Les deux corrections vont en sens opposés et se compensent presque exactement sur la
\emph{perte moyenne} ($0{,}861\times1{,}147=0{,}987$). Elles ne se compensent \textbf{pas} sur
le capital : à queue lourde, le quantile suit la sévérité et non la fréquence. Deux biais qui
s'annulent sur l'espérance peuvent ne pas s'annuler sur le quantile, et c'est le quantile qui
fait le SCR.}

% =====================================================================
\part*{II. Combien d'incidents par an}
\addcontentsline{toc}{part}{II. Combien d'incidents par an}

\section{Le processus de Poisson, et pourquoi il ne suffit pas}

\lequoi{Modéliser un nombre d'événements par an, sans savoir quand ils tombent.}

\lidee{Si les événements arrivent indépendamment, à un rythme constant, et jamais deux
exactement en même temps, alors leur nombre annuel suit une loi de Poisson. C'est le modèle
de comptage par défaut, et il n'a qu'un paramètre.}

\laformule{$\mathbb{P}(N=k)=e^{-\lambda}\lambda^k/k!$, avec la propriété caractéristique
\[
  \mathbb{E}[N]=\mathrm{Var}(N)=\lambda .
\]
Moyenne et variance sont \emph{égales}. C'est à la fois sa force et sa faiblesse.}

\cheznous{Cette égalité est testable, et elle est rejetée massivement. Sur $14\,944$
cellules entité-année, un test du rapport de vraisemblance donne $p\approx10^{-30}$.
Autrement dit~: la probabilité d'observer nos données si Poisson était vrai est de l'ordre
d'un sur $10^{30}$.}

\laphrase{Poisson impose que la variance égale la moyenne. Dans les données, la variance est
neuf fois la moyenne. Ce n'est pas une nuance, c'est un rejet à $10^{-30}$, donc j'ai changé
de loi.}

\lepiege{Ne pas dire \og Poisson sous-estime le risque\fg{}. Poisson donne la \emph{bonne}
moyenne, il se trompe sur la \emph{dispersion}, donc sur la queue. Le SCR est un quantile
extrême~: c'est la dispersion qui le fait, pas la moyenne.}

\section{La surdispersion et la négative binomiale}

\lequoi{Les données montrent beaucoup plus de variabilité que Poisson n'en autorise. Pourquoi,
et comment en tenir compte ?}

\lidee{Parce que le rythme lui-même varie. Certaines années sont plus exposées que
d'autres, certaines entités plus que d'autres. On rend donc $\lambda$ \emph{aléatoire} au
lieu de le fixer. Un mélange de Poisson par une loi Gamma donne exactement la négative
binomiale.}

\laformule{Si $N\mid\Lambda\sim\mathcal{P}(\Lambda)$ et $\Lambda\sim\mathrm{Gamma}$, alors
$N$ est négative binomiale et
\[
  \mathbb{E}[N]=\lambda, \qquad \mathrm{Var}(N)=\varphi\,\lambda, \qquad \varphi>1 .
\]
Le facteur $\varphi$ est la \textbf{surdispersion}~: le rapport variance sur moyenne.}

\cheznous{$\varphi=9{,}2$. La variance vaut neuf fois la moyenne. Trois mesures de cette
surdispersion coexistaient dans le modèle sans avoir jamais été confrontées~: deux d'entre
elles se sont révélées être la même quantité vue de deux angles, elles coïncident à
$5\,\%$. À variance égale, la \emph{forme} du mélange (Gamma ou lognormale) déplace encore
la queue de $11\,\%$.}

\laphrase{La négative binomiale, c'est Poisson dont le rythme est lui-même incertain.
Concrètement, certaines années sont plus exposées que d'autres, et ce sont ces années-là
qui font le capital.}

\lepiege{Le changement de loi de comptage \textbf{ne déplace pas la moyenne}, il déplace la
VaR. C'est logique~: $\mathbb{E}[N]=\lambda$ dans les trois lois testées. Sur la série
détrendée, la surdispersion est \emph{trois fois plus faible} que celle retenue~: le calage
est conservateur, pas calibré, et il faut le dire soi-même.}

\section{Pourquoi pas un processus de Hawkes}

\lequoi{Un processus auto-excité, où chaque événement augmente temporairement la
probabilité du suivant, serait un candidat naturel pour de la contagion.}

\lidee{Hawkes modélise la contagion \emph{dans le temps}~: un incident aujourd'hui rend un
autre incident plus probable demain. Notre cascade modélise la contagion \emph{entre
domaines}~: un incident sur un pilier rend un incident plus probable sur un autre pilier.
Ce ne sont pas les mêmes axes.}

\cheznous{L'estimation d'un noyau d'excitation demande des dates fines. Nos données sont
annuelles~: le noyau serait entièrement non identifié, exactement comme la direction. La
littérature cyber sur Hawkes (Boumezoued, Hillairet) travaille sur des données horodatées
finement, que nous n'avons pas.}

\laphrase{Hawkes décrit la contagion dans le temps, ma cascade décrit la contagion entre
domaines. Les deux sont compatibles, mais Hawkes demande des dates fines que mes données
n'ont pas. Je ne le contredis pas, je ne peux simplement pas le calibrer.}

\lepiege{Ne pas présenter le rejet de Hawkes comme une critique de cette littérature. La
formulation qui marche est \og cela ne contredit pas, cela ne s'applique pas à ma donnée\fg{}.}

% =====================================================================
\part*{III. Combien coûte un incident}
\addcontentsline{toc}{part}{III. Combien coûte un incident}

\section{Pourquoi la théorie des valeurs extrêmes}

\lequoi{Le SCR est un quantile à $99{,}5\,\%$. On veut donc décrire précisément la région
où l'on a le moins d'observations. Ajuster une loi sur l'ensemble des données optimise
l'ajustement là où il y a le plus de points, c'est-à-dire au \emph{centre}, exactement là
où l'on s'en moque.}

\lidee{On abandonne l'idée d'une loi unique pour tout. On modélise \emph{seulement} la
queue, avec une famille dont on peut démontrer qu'elle est la seule limite possible.}

\laphrase{Une loi ajustée sur toutes les données est calée sur le centre, or je n'ai besoin
que de la queue. La théorie des valeurs extrêmes dit quelle famille de lois est la seule
possible tout au bout, et je n'ajuste que celle-là.}

\section{Les deux théorèmes fondateurs}

\lidee{Deux résultats, structurés exactement comme le théorème central limite~: ils disent
qu'\emph{une seule} famille peut apparaître à la limite, quelle que soit la loi de départ.}

\laformule{\textbf{Fisher-Tippett-Gnedenko.} Si le maximum renormalisé de $n$ variables
indépendantes converge en loi, la limite ne peut être que la loi des valeurs extrêmes
généralisée~:
\[
  G_\xi(x)=\exp\!\left\{-\left(1+\xi x\right)^{-1/\xi}\right\}, \qquad 1+\xi x>0 .
\]
\textbf{Pickands-Balkema-de Haan.} De façon équivalente, la loi des \emph{excès} au-dessus
d'un seuil $u$ converge, quand $u$ monte vers le point terminal, vers la loi de Pareto
généralisée~:
\[
  \mathbb{P}(X-u\le y \mid X>u)\;\longrightarrow\;
  1-\left(1+\frac{\xi y}{\sigma}\right)^{-1/\xi}.
\]}

\cheznous{C'est le second qu'on utilise, parce qu'il exploite \emph{tous} les dépassements
de seuil et pas seulement le maximum de chaque bloc. On perd beaucoup moins d'information.}

\laphrase{C'est le même genre de résultat que le théorème central limite. Là où le TCL dit
que les moyennes tendent vers une gaussienne quelle que soit la loi de départ, ces deux
théorèmes disent que les extrêmes tendent vers une seule famille. Je n'ai donc pas
\emph{choisi} la loi de queue, elle est imposée.}

\lepiege{Ces théorèmes sont \emph{asymptotiques}. Ils ne garantissent pas que la GPD est
correcte à un seuil fini~: ils garantissent qu'aucune autre famille ne peut l'être à la
limite. C'est plus faible, il faut le dire, et c'est pourquoi on teste l'adéquation.}

\section{La méthode des excès et la loi de Pareto généralisée}

\lidee{On fixe un seuil $u$. En dessous, on garde la distribution empirique telle quelle. Au
dessus, on ajuste une GPD. La loi obtenue est dite \emph{recollée} (spliced).}

\laformule{La fonction de survie s'écrit, pour $x>u$~:
\[
  \mathbb{P}(X>x)\;=\;p_u\left(1+\frac{\xi (x-u)}{\sigma}\right)^{-1/\xi},
  \qquad p_u=\mathbb{P}(X>u).
\]}

\cheznous{$\xi=0{,}60$ sur OpRisk, $\xi=1{,}03$ sur PRC. Le paramètre $\xi$ est
l'\textbf{indice de queue}, et il commande tout~:
\begin{itemize}[itemsep=1pt,topsep=2pt]
\item $\xi<1/2$~: variance finie ;
\item $1/2<\xi<1$~: variance \emph{infinie}, moyenne finie ;
\item $\xi>1$~: moyenne elle-même infinie.
\end{itemize}
Notre $\xi=1{,}03$ sur PRC dépasse $1$~: la moyenne théorique n'existe pas. Ce n'est pas une
erreur de calcul, c'est une propriété du cyber.}

\laphrase{L'indice de queue mesure à quel point les catastrophes dominent. Au dessus de un,
la moyenne théorique n'existe même plus~: c'est un domaine où un seul sinistre peut peser
plus que tous les autres réunis. C'est exactement ce qu'on observe en cyber.}

\lepiege{Avec $\xi$ proche de $1$, la moyenne empirique n'est pas un estimateur stable et
les intervalles de confiance sont larges par nature. Ne jamais annoncer un SCR ponctuel
sans sa bande dans ce régime.}

\section{Le choix du seuil}

\lequoi{Le seuil $u$ arbitre entre biais et variance~: trop bas, on inclut le corps de la
distribution et l'approximation GPD est fausse ; trop haut, il reste trop peu de points et
l'estimation devient bruitée.}

\lidee{On ne le choisit pas au hasard, on le lit sur deux graphiques de diagnostic.}

\laformule{\textbf{Vie résiduelle moyenne.} Pour une GPD, la fonction
$e(u)=\mathbb{E}[X-u\mid X>u]$ est \emph{linéaire} en $u$. On choisit $u$ au début de la
zone linéaire.
\textbf{Graphique de Hill.} On trace $\hat\xi$ en fonction de $u$ et on cherche un plateau.}

\cheznous{Seuil retenu $u=4{,}176$~M\euro{} sur PRC, avec $p_u=0{,}150$, donc $15\,\%$ des
pertes dépassent le seuil. Les deux diagnostics concordent.}

\laphrase{Le seuil se lit sur deux graphiques~: la vie résiduelle moyenne doit devenir
droite, et l'indice de queue doit se stabiliser. Les deux disent la même chose ici, et c'est
ce qui rend le choix défendable plutôt qu'arbitraire.}

% =====================================================================
\part*{IV. Comment les pertes s'additionnent}
\addcontentsline{toc}{part}{IV. Comment les pertes s'additionnent}

\section{Le modèle collectif}

\lidee{On ne modélise pas la perte annuelle directement. On modélise séparément
\emph{combien} d'incidents et \emph{combien coûte} chacun, puis on somme. C'est le cadre
standard en actuariat non-vie.}

\laformule{
\[
  L \;=\; \sum_{i=1}^{N} X_i,
\]
avec $N$ le nombre d'incidents (fréquence) et $X_i$ leurs coûts (sévérité), les $X_i$
indépendants de $N$ et entre eux.}

\cheznous{$10^5$ années simulées par Monte-Carlo. On ne cherche pas de formule fermée pour
la loi de $L$~: avec une queue lourde et une cascade, il n'y en a pas.}

\laphrase{Je sépare deux questions qui n'ont rien à voir~: combien d'incidents par an, et
combien coûte un incident. Chacune se calibre sur ses propres données, et la perte annuelle
est leur somme.}

\lepiege{L'hypothèse d'indépendance entre $N$ et les $X_i$ est une hypothèse, pas un fait.
Dans une crise, il y a plus d'incidents \emph{et} ils sont plus chers. C'est exactement ce
que la cascade réintroduit.}

\section{VaR, TVaR, et pourquoi le SCR est une VaR}

\lidee{La VaR à $99{,}5\,\%$ est le montant qu'on ne dépasse que dans une année sur deux
cents. La TVaR est la moyenne des pertes \emph{au-delà} de ce seuil.}

\laformule{
\[
  \mathrm{VaR}_\alpha(L)=\inf\{x:\ \mathbb{P}(L\le x)\ge\alpha\},
  \qquad
  \mathrm{TVaR}_\alpha(L)=\mathbb{E}[L\mid L>\mathrm{VaR}_\alpha(L)].
\]}

\cheznous{Solvabilité II impose la VaR à $99{,}5\,\%$ à horizon un an. On ne choisit pas
cette mesure, elle est réglementaire, et c'est un point à faire valoir quand on nous
reproche d'utiliser un quantile.}

\laphrase{Le SCR, c'est la perte qu'on ne dépasse qu'une année sur deux cents. Ce n'est pas
mon choix de mesure, c'est celui de Solvabilité II.}

\lepiege{La VaR n'est pas sous-additive~: la VaR d'une somme peut dépasser la somme des VaR.
La TVaR, elle, est cohérente. Si on nous interroge sur ce point~: \og je rapporte les deux, et
la TVaR est effectivement la mesure la plus saine, mais le capital réglementaire se calcule
en VaR\fg{}.}

% =====================================================================
\part*{V. Comment un problème en entraîne un autre}
\addcontentsline{toc}{part}{V. Comment un problème en entraîne un autre}

\section{La convention d'indices, à ne jamais improviser}

\lequoi{Une matrice dirigée a deux lectures possibles, et les deux sont défendables. Si le
document en utilise l'une ici et l'autre là, tous ses énoncés directionnels deviennent
ambigus.}

\lidee{On fixe la convention une fois, on l'écrit, et on la vérifie sur un fait connu.}

\laformule{\[
  W_{jk}\;=\;\text{propension de } j \text{ (\textbf{source}) à entraîner la défaillance de }
  k \text{ (\textbf{cible}).}
\]
La \textbf{ligne} $j$ décrit ce que $j$ \emph{émet}, la \textbf{colonne} $k$ ce que $k$
\emph{reçoit}. La somme de ligne $s_j$ mesure donc une \textbf{émission}.}

\cheznous{Vérification sur la matrice d'expert : P1 a la plus grande somme de ligne ($2{,}6$)
et une somme de colonne trois fois inférieure à celle de P2. C'est exactement la revendication
centrale \og{}P1 émet fort et reçoit peu\fg{}, et elle serait \emph{fausse} sous la convention
inverse. La forme réduite du modèle latent s'écrit alors $(I-W^{\top})^{-1}$, le rayon spectral
étant de toute façon invariant par transposition.}

\laphrase{La ligne émet, la colonne reçoit. Je le vérifie sur ma matrice : le pilier que je dis
émetteur a bien la plus grosse somme de ligne. Une convention n'est pas un détail de notation,
c'est ce qui rend un énoncé directionnel falsifiable.}

\lepiege{Ce défaut a réellement existé dans le mémoire : l'équation structurelle sommait sur
les lignes alors que la définition de la cascade et tout le code sommaient sur les colonnes. Le
code et le classeur ont tranché. Vérifier une convention sur un fait connu coûte trente
secondes et évite une matrice transposée.}

\section{La matrice de contagion et sa normalisation}

\lequoi{On veut dire \og un problème sur le pilier $j$ aggrave le pilier $k$\fg{}, avec une
intensité. Mais si on met des coefficients au hasard, la cascade peut exploser.}

\begin{center}
\begin{tikzpicture}[scale=1.0]
\node[pilier] (p1) at (90:2.0)  {P1};
\node[pilier] (p2) at (18:2.0)  {P2};
\node[pilier] (p3) at (-54:2.0) {P3};
\node[pilier] (p4) at (234:2.0) {P4};
\node[pilier] (p5) at (162:2.0) {P5};
\node[etiq,above=1pt of p1] {gouvernance TIC};
\node[etiq,above right=0pt and 1pt of p2] {incidents};
\node[etiq,below=1pt of p3] {tests};
\node[etiq,below=1pt of p4] {tiers};
\node[etiq,left=1pt of p5]  {partage};
% P1 emet fort, recoit peu
\draw[flechef] (p1) to[bend left=12] (p2);
\draw[flechef] (p1) to[bend left=12] (p3);
\draw[flechef] (p1) to[bend right=12] (p4);
\draw[fleche]  (p4) to[bend right=12] (p2);
\draw[fleche]  (p4) to[bend left=12]  (p3);
\draw[flechep] (p5) to[bend left=12]  (p1);
\draw[flechep] (p3) to[bend left=12]  (p2);
\node[align=left,font=\scriptsize,text=navy!75,anchor=west] at (4.3,0.9)
  {\textbf{Lire la matrice~:}\\$W_{jk}$ = propension de $j$\\à entraîner $k$.\\
   \emph{La ligne émet, la colonne reçoit.}};
\node[align=left,font=\scriptsize,text=accent,anchor=west] at (4.3,-0.9)
  {P1 émet fort et reçoit peu.\\C'est pourquoi il sort\\premier en vue Shapley.};
\end{tikzpicture}
\end{center}

\noindent La cascade est un \emph{processus de branchement}~: un incident amorcé sur un pilier
tire, pour chaque arc, une pièce de Bernoulli de probabilité $W_{jk}$, et recommence depuis les
piliers atteints. Sans garde-fou, ce processus peut ne jamais s'éteindre.

\lidee{On part du jugement dirigé (la matrice \textsc{trans} du classeur), et on la
renormalise pour garantir mathématiquement que la propagation s'éteint.}

\laformule{
\[
  W \;=\; g\,\frac{\mathrm{TRANS}}{\max_j \sum_k \mathrm{TRANS}_{jk}},
  \qquad g=0{,}90 ,
\]
ce qui garantit $\rho(W)\le g<1$, où $\rho$ est le rayon spectral (la plus grande valeur
propre en module).}

\cheznous{Le facteur $g$ est un curseur de propagation~: $g=0$ signifie aucune contagion,
$g\to1$ signifie une cascade au bord de l'explosion. On le fixe à $0{,}90$, ce qui est
volontairement élevé.}

\laphrase{La normalisation garantit que la cascade s'éteint. Sans elle, un choc pourrait
s'amplifier indéfiniment en tournant en boucle entre les piliers, et le capital serait
infini.}

\lepiege{$g=0{,}90$ est un \emph{choix}, pas une mesure. À dire ainsi~: \og c'est un
paramètre de sévérité de scénario, et je montre le résultat pour $g=0$ et $g=0{,}90$ afin
qu'on voie ce qu'il porte\fg{}.}

\section{L'inverse de Leontief et la sous-criticité}

\lidee{Le choc initial se propage, ce qu'il déclenche se propage à son tour, et ainsi de
suite. La somme de tous les tours est une série géométrique matricielle.}

\laformule{
\[
  (I-W)^{-1} \;=\; I + W + W^2 + W^3 + \cdots
\]
Cette série converge \emph{si et seulement si} $\rho(W)<1$. Le terme $I$ est le choc
initial, $W$ ce qu'il déclenche directement, $W^2$ ce que cela déclenche à son tour.}

\cheznous{L'objet vient de l'économie (Leontief, tableaux entrées-sorties)~: la même algèbre
décrit comment une commande dans un secteur en génère dans les autres.}

\laphrase{C'est la même formule que les tableaux entrées-sorties en économie. Le premier
terme est le choc direct, le deuxième ce qu'il déclenche, le troisième ce que cela déclenche
à son tour. La condition de convergence est exactement la condition pour que la cascade
s'éteigne.}

\lepiege{Le passage à $W^2$ est celui où la \emph{direction} commence à compter, alors
qu'elle ne compte pas dans $W$ seul si on ne regarde que les co-occurrences. C'est le point
technique le plus important du mémoire, voir la section~\ref{sec:trace}.}

\section{Le seuil de déclenchement, et le rôle de la loi normale}

\lequoi{À partir de quand un pilier bascule ? Il faut un mécanisme, pas seulement une
probabilité.}

\lidee{On suppose une variable latente $X_j$ qui mesure la fragilité du pilier $j$, et le
pilier bascule quand elle franchit un seuil. C'est le mécanisme de Merton, repris par
Vasicek pour le risque de crédit.}

\laformule{Le pilier $j$ bascule quand
\[
  X_j \;\ge\; K_j \;=\; F^{-1}(1-p_j),
\]
où $p_j$ est sa probabilité marginale de basculement et $F$ la fonction de répartition de la
latente.}

\cheznous{$K_4 = 1{,}033$ pour le pilier prestataires.}

\laphrase{Il y a deux questions derrière le mot seuil, et il ne faut pas les mélanger. Pour
savoir \emph{quand un pilier bascule tout seul}, la loi normale est inoffensive~: elle ne
sert que d'unité de mesure, la vraie queue lourde est dans la sévérité, pas ici. Pour savoir
\emph{combien de piliers basculent ensemble}, c'est autre chose, et là le gaussien pose un
vrai problème.}

\lepiege{Le choix de la loi de la latente est une \emph{normalisation} pour la marge~:
n'importe quelle loi continue donnerait la même probabilité marginale, seul le seuil change
de valeur. Ce n'est pas vrai pour la dépendance. Confondre les deux est l'erreur classique.}

\section{Les copules et la dépendance de queue}

\lequoi{Comment plusieurs piliers basculent-ils \emph{ensemble} ? La probabilité marginale
de chacun ne suffit pas à le dire.}

\lidee{Une copule sépare les marges de la structure de dépendance. On peut garder les mêmes
marges et changer complètement la façon dont les événements extrêmes se coordonnent.}

\laformule{Le coefficient de \textbf{dépendance de queue supérieure} mesure la
co-occurrence extrême~:
\[
  \lambda_U \;=\; \lim_{q\to1^-}\ \mathbb{P}\bigl(U_2>q \mid U_1>q\bigr).
\]
Pour la copule gaussienne, $\lambda_U = 0$ dès que la corrélation est $<1$. Pour la copule
de Gumbel de paramètre $\theta$,
\[
  \lambda_U \;=\; 2-2^{1/\theta}.
\]}

\cheznous{Avec $\theta=1{,}8$, on a $\lambda_U = 0{,}53$. À $\tau$ de Kendall \emph{égal},
la gaussienne donne $0$ et Gumbel $0{,}53$. La probabilité que les cinq piliers basculent
ensemble passe de $0{,}17$ à $0{,}39$ selon la copule retenue. On encadre au lieu de
trancher~: plancher gaussien, plafond de Gumbel.}

\laphrase{Une dépendance gaussienne a une propriété gênante~: dans les extrêmes, elle prédit
que les choses se \emph{découplent}. C'est le contraire de ce qu'on observe en crise. Donc
je n'ai pas tranché entre les deux, j'ai encadré, et l'écart entre le plancher et le plafond
est un résultat en soi.}

\lepiege{Deux copules peuvent avoir le \emph{même} $\tau$ de Kendall, donc la même
dépendance \og moyenne\fg{}, et des comportements extrêmes opposés. Comparer deux copules par
leur corrélation ne dit rien sur la queue.}

% =====================================================================
\part*{VI. Ce qu'on ne peut pas savoir, et pourquoi}
\addcontentsline{toc}{part}{VI. Ce qu'on ne peut pas savoir, et pourquoi}

\section{La décomposition $W = S + A$}

\lidee{Toute matrice se coupe de manière unique en une partie symétrique et une partie
antisymétrique. La première dit \emph{qui est lié à qui}, la seconde dit \emph{dans quel
sens}.}

\laformule{
\[
  S=\tfrac12\bigl(W+W^{\top}\bigr), \qquad A=\tfrac12\bigl(W-W^{\top}\bigr),
  \qquad W=S+A .
\]
Les deux parties sont \textbf{orthogonales}~: $\langle S,A\rangle=0$, d'où un Pythagore
$\|W\|^2=\|S\|^2+\|A\|^2$.}

\cheznous{$\|A\|/\|W\| = 40{,}7\,\%$. La direction représente 40\,\% de la matrice, ce n'est
pas un détail.}

\laphrase{On coupe la matrice en deux~: une partie qui dit quels piliers sont liés, et une
partie qui dit dans quel sens. Les deux sont orthogonales, ce qui veut dire que connaître
l'intensité n'apprend rien sur la direction.}

\lepiege{La décomposition propre porte sur le \emph{flux stationnaire} $\pi_j P_{jk}$, pas
sur la matrice de transition brute. Notre $S=(W+W^{\top})/2$ en est le cas où $\pi$ est
uniforme, et l'écart mesuré est de $0{,}069$, donc faible. Le dire soi-même met à l'abri.}

\medskip
\noindent\textbf{Le schéma qui fait comprendre la thèse en dix secondes.} Deux matrices
opposées, une seule statistique observable, deux décisions contraires.

\begin{center}
\begin{tikzpicture}
% cas W
\node[pilier] (a1) at (0,1.1) {P1};
\node[pilier] (b1) at (0,-0.5) {P2};
\draw[flechef] (a1) -- (b1);
\node[font=\scriptsize\bfseries,text=navy] at (0,1.9) {matrice $W$};
\node[etiq,align=center] at (0,-1.3) {P1 entraîne P2\\\emph{remédier P1}};
% cas W transpose
\node[pilier] (a2) at (3.4,1.1) {P1};
\node[pilier] (b2) at (3.4,-0.5) {P2};
\draw[flechef] (b2) -- (a2);
\node[font=\scriptsize\bfseries,text=navy] at (3.4,1.9) {matrice $W^{\top}$};
\node[etiq,align=center] at (3.4,-1.3) {P2 entraîne P1\\\emph{remédier P2}};
% ce que la donnee voit
\node[draw=vert!70!black,fill=vert!8,rounded corners=2pt,align=center,
      font=\footnotesize,inner sep=6pt] (obs) at (8.4,0.3)
  {\textbf{Ce que la donnée agrégée voit}\\
   \og P1 et P2 tombent ensemble\fg{}\\[2pt]
   $\Prob(j \text{ et } k) $ ne dépend que de $S$\\
   et $S(W)=S(W^{\top})$};
\draw[fleche] (1.0,0.3) -- (obs.west |- 0,0.3);
\draw[fleche] (4.4,0.3) -- (obs.west |- 0,0.3);
\node[align=center,font=\scriptsize,text=accent] at (4.25,-2.3)
  {\textbf{Indiscernables pour la donnée. Matériellement différentes pour la décision.}};
\end{tikzpicture}
\end{center}

\noindent C'est l'énoncé exact de la frontière d'identification, et il ne repose sur
\emph{aucun} dire d'expert~: la co-occurrence est une fonction de $S$ seule, et $S$ est
identique pour $W$ et pour sa transposée. Toute l'information directionnelle vit dans $A$, que
la co-occurrence ne voit pas.

\section{Le renversement du temps et la fluctuation martingale}

\lequoi{Pourquoi n'arrive-t-on pas à mesurer $A$ ? Et est-ce une limite de méthode ou une
limite de principe ?}

\lidee{Passer le film à l'envers. La partie symétrique ne change pas, la partie
antisymétrique s'inverse. Et la fluctuation d'un processus, une fois la tendance retranchée,
ne dépend que de la première.}

\laformule{La chaîne renversée a pour flux $F^{*}=F^{\top}=S-A$. Et pour toute fonction test
$f$, l'intensité de la fluctuation vaut
\[
  \mathbb{E}_\pi\bigl[\Gamma(f,f)\bigr]
  \;=\;\sum_{j,k} S_{jk}\bigl(f(k)-f(j)\bigr)^2 ,
\]
la partie $A$ y contribuant exactement zéro. La démonstration tient en une ligne~:
$(f(k)-f(j))^2$ est symétrique en $(j,k)$, $A$ est antisymétrique, les termes $(j,k)$ et
$(k,j)$ se compensent.}

\cheznous{Vérifié sur $2\,000$ fonctions test tirées au hasard~: écart maximal
$1{,}8\times10^{-15}$, c'est-à-dire zéro numérique.}

\laphrase{Renverser le temps laisse la partie symétrique intacte et retourne le signe de la
partie directionnelle. Or la fluctuation d'un processus, sa partie martingale, ne dépend que
de la symétrique. Donc une donnée annuelle, qui ne garde pas l'ordre, ne peut
mathématiquement pas voir la direction.}

\lepiege{Ne pas dire que $W$ \emph{est} une chaîne de Markov. C'est un noyau de contagion~:
on emprunte l'algèbre de la réversibilité comme cadre de lecture, on n'ajoute pas une
hypothèse.}

\section{Le test placebo}
\label{sec:placebo}

\lequoi{Même des données sans aucune direction produisent un déséquilibre, par simple
fluctuation d'échantillonnage. Observer un déséquilibre ne prouve donc rien tant qu'on ne
sait pas ce que le hasard produit tout seul.}

\lidee{On fabrique une version des données dont la direction a été \textbf{détruite
volontairement}, tout le reste étant conservé, et on compare. C'est l'essai clinique~: si le
médicament ne fait pas mieux que la pilule vide, il ne fait rien.}

\laformule{On mesure l'asymétrie de la matrice de transitions comptées~:
\[
  \mathcal{A}(M)=\frac{\|M-M^{\top}\|_1}{2}
  =\frac12\sum_{j,k}\bigl|M_{jk}-M_{kj}\bigr| ,
\]
c'est-à-dire le surplus de transitions d'un sens sur l'autre, sommé sur toutes les paires.
Puis on \textbf{permute les années à l'intérieur de chaque entité}~: l'ensemble des
catégories touchées chaque année est conservé, seul l'ordre est détruit. On recompte, on
recalcule, et on compare~:
\[
  z=\frac{\mathcal{A}(M_{\text{obs}})-\overline{\mathcal{A}^{\,\text{null}}}}
        {\widehat{s}(\mathcal{A}^{\,\text{null}})} .
\]}

\cheznous{
\begin{center}\small
\begin{tabular}{@{}l r r r l@{}}
\toprule
\textbf{Source} & \textbf{observée} & \textbf{sous placebo} & \textbf{$z$} & \textbf{Verdict}\\
\midrule
OpRisk, pas annuel & $119$ & $128\pm27$ & $-0{,}33$ & aucun signal \\
7 post-mortems & \multicolumn{2}{c}{\emph{même protocole}} & $+3{,}93$ & signal net \\
\bottomrule
\end{tabular}
\end{center}
Sur OpRisk, l'observé est \emph{sous} la moyenne du bruit. Un test binomial par paire le
confirme~: aucune paire ne survit à la correction de Bonferroni.}

\laphrase{Je voulais savoir si la direction que je crois voir est réelle. Le problème, c'est
que même des données au hasard produisent un déséquilibre. Alors j'ai pris mes données, j'ai
détruit l'ordre en gardant tout le reste, et j'ai comparé. Mes vraies données font
$119$, le hasard fait $128$~: elles font moins bien que le hasard. Puis j'ai refait
exactement le même test sur les sept rapports d'enquête, et là le signal apparaît nettement.
Même test, autre source, résultat opposé~: le problème venait de la donnée, pas de la
méthode.}

\lepiege{Trois choses à ne pas dire. \textbf{(i)} Le test du \og temps inversé\fg{} est
\emph{dégénéré} sur un comptage~: renverser la chronologie donne $M^{\top}$ par
construction, donc une corrélation de $-1$ quel que soit le contenu. Il ne prouve rien.
\textbf{(ii)} Un test non rejeté n'est pas une preuve d'absence~: la donnée est
\emph{compatible} avec l'absence de direction. \textbf{(iii)} Sur les post-mortems, le test
établit la \emph{cohérence} du codage, pas son \emph{exogénéité}~: le biais de narration des
rapports d'enquête n'est pas levé par le calcul.}

\section{La production d'entropie}

\lidee{Un seul nombre qui dit à quel point le processus est dirigé, et qui a une
interprétation exacte~: à quel point le film à l'endroit est distinguable du film à
l'envers.}

\laformule{
\[
  \sigma=\sum_{j<k} 2A_{jk}\,\ln\!\frac{S_{jk}+A_{jk}}{S_{jk}-A_{jk}} \;\ge\; 0,
\]
nulle si et seulement si $A=0$. Elle est aussi le taux d'entropie relative entre la loi de
la trajectoire lue à l'endroit et celle lue à l'envers.}

\cheznous{$\sigma=0{,}0700$ pour la chaîne posée, $5\times10^{-33}$ après symétrisation.
\textbf{Et le résultat qui compte}~: chaque terme de la somme est \emph{pair} en $A_{jk}$,
donc $\sigma$ ne dépend que des \emph{amplitudes}, pas des signes. Vérifié numériquement~:
les $1024$ sommets de l'ensemble admissible ont tous \emph{exactement} la même $\sigma$,
amplitude $0{,}0\times10^{0}$, alors que le capital y varie de $27\,\%$.}

\laphrase{La production d'entropie mesure à quel point on peut distinguer le film à
l'endroit du film à l'envers. Si elle est nulle, personne ne peut dire dans quel sens le
temps coule. Mais attention~: elle ne voit que la \emph{quantité} de direction, jamais son
\emph{sens}. Les mille vingt-quatre configurations ont toutes la même production d'entropie,
et pourtant elles étalent le capital d'un quart.}

\lepiege{Conséquence à tirer~: ce n'est pas $\sigma$ qu'il faut chercher à estimer, ce sont
les \emph{signes}, un par un. C'est ce qui justifie que la valeur d'information se
hiérarchise par paire et non globalement.}

\section{$z$ scalaire ou $z$ par paire}

\lequoi{L'écriture $P=S+zA$ est ambiguë tant qu'on n'a pas dit si $z$ est un nombre ou dix.}

\laformule{\textbf{Lecture globale}~: $z\in[-1,1]$ règle l'\emph{intensité} d'un motif de
direction fixé. Un seul degré de liberté. Elle ne peut pas changer quel pilier précède
quel autre.
\textbf{Lecture par paire}, celle du mémoire~:
\[
  A_{jk}=z_{jk}\,t\,S_{jk}, \qquad z_{jk}\in[-1,1], \qquad j<k,
\]
soit \textbf{dix degrés de liberté}. Le paramètre $t$ est l'amplitude maximale d'asymétrie,
les $z_{jk}$ portent chacun un sens indépendant.}

\cheznous{Les $2^{10}=1024$ sommets sont exactement les configurations $z_{jk}=\pm1$. La
contrainte de positivité du flux impose $|A_{jk}|\le S_{jk}$, donc $t\le1$~; $t=1$ est
l'ignorance totale.}

\laphrase{Un $z$ global ne fait que renforcer ou affaiblir un ordre déjà choisi. Ce qu'il me
faut, c'est un curseur par paire, parce que c'est le \emph{sens} de chaque lien que je ne
connais pas, pas son intensité globale.}

\section{Le troisième échec : le schéma existe, il n'est pas rempli}

\lequoi{Deux bases ont échoué à identifier la direction, l'une faute de taxonomie, l'autre
faute de résolution temporelle. On pouvait encore croire que personne n'avait \emph{prévu} de
collecter la bonne donnée.}

\lidee{On teste l'hypothèse la plus favorable : un standard de description d'incidents qui,
lui, prévoit exactement les bons champs. Le schéma \textsc{veris}, référence de la profession
et support du \emph{Data Breach Investigations Report}, porte littéralement les deux objets
manquants : un champ \texttt{control\_failure} qui nomme le domaine de contrôle défaillant, et
une chronologie en quatre étapes qui porte l'antériorité.}

\cheznous{Sur les $10\,591$ incidents de la base publique associée : \texttt{control\_failure}
renseigné dans \textbf{$0{,}12\,\%$} des cas ($13$ incidents, dont un seul en secteur
financier), chronologie datée sur au moins deux étapes dans \textbf{$1{,}19\,\%$}. Et
l'\textbf{intersection} des deux, qui est ce qu'identifier $W$ exigerait simultanément, ne
compte qu'\textbf{un seul incident}, et \textbf{aucun} en secteur financier.}

\laphrase{Le champ qui donnerait la réponse existe dans le standard de la profession. Il est
rempli une fois sur mille. Le problème n'est donc pas de savoir quoi collecter, il est d'en
rendre le remplissage obligatoire, et c'est précisément ce qu'un régime de notification
réglementaire peut imposer là où une base déclarative volontaire ne le peut pas.}

\lepiege{C'est le plus utile des trois échecs, parce qu'il transforme la recommandation. On ne
demande plus d'inventer une taxonomie, ce qui serait long et contestable, on demande de rendre
obligatoires deux champs \textbf{déjà normalisés}. Coût de conception nul.}

\section{La matrice ordonnée $p_{jk}$}

\lequoi{Le codage documentaire, tel qu'on l'exploite d'abord, compare $n(j\to k)$ à
$n(k\to j)$ \emph{à l'intérieur} d'une paire. Ce rapport élimine le dénominateur, donc il
donne un \emph{sens} et ne dit rien de l'\emph{amplitude}.}

\lidee{On ajoute le dénominateur. Chaque incident déclare non seulement ses transitions, mais
l'ensemble des \textbf{piliers défaillants}, y compris ceux qui ont cédé sans rien entraîner.
On peut alors estimer directement chaque case de la matrice ordonnée.}

\laformule{Avec $N_j$ le nombre d'incidents où $j$ a cédé et $m_{jk}$ ceux où le rapport
établit $j\to k$ :
\[
  p_{jk}=\Prob(k \text{ cède}\mid j \text{ a cédé}),
  \qquad
  p_{jk}\mid\text{données}\ \sim\ \mathrm{Beta}\bigl(1+m_{jk},\ 1+N_j-m_{jk}\bigr).
\]
C'est la matrice ordonnée \emph{complète}, donc $S$ \textbf{et} $A$.}

\cheznous{Sur dix rapports et $22$ transitions : \textbf{douze} liens sur vingt ont un crédible
à $90\,\%$ entièrement \emph{sous} $0{,}5$, et \textbf{aucun} au-dessus. Les plus élevés,
$P_3\to P_2$ ($0{,}77$, $[0{,}48\,;0{,}95]$) et $P_1\to P_4$ ($0{,}68$, $[0{,}40\,;0{,}89]$),
recouvrent encore $0{,}5$. Les douze liens faibles comprennent les cinq sorties de $P_2$ et les
quatre entrées de $P_1$.}

\laphrase{Mon corpus établit très bien où la contagion \emph{ne va pas}, et pas du tout où elle
va fort. Le puits et la source pure sont donc établis en tant que tels, mais aucun coefficient
n'est établi comme grand. C'est un ordre, ce n'est pas une magnitude.}

\lepiege{Une enquête raconte la \emph{chaîne} qui a produit le sinistre : elle documente bien
les contrôles qui ont cédé \emph{et} propagé, mal ceux qui ont cédé sans propager, pas du tout
ceux qui ont tenu. Le dénominateur $N_j$ est donc sous-estimé et $p_{jk}$ est une \textbf{borne
supérieure}. En ajoutant des défaillances non propagées, le \emph{niveau} s'effondre de
$0{,}27$ à $0{,}07$ tandis que la part de direction $\lVert A\rVert/\lVert W\rVert$ ne bouge
pas ($60\,\%$). Le corpus établit un ordre, pas une magnitude, et cette phrase-là est robuste.}

\section{L'identification partielle}

\lidee{Quand un paramètre n'est pas identifiable, il reste deux attitudes~: le poser
arbitrairement et faire semblant, ou énumérer tout ce qu'il pourrait être et rapporter
l'intervalle qui en résulte. La seconde est l'identification partielle, au sens de Manski.}

\laformule{On définit l'ensemble admissible
\[
  \mathcal{W}=\bigl\{\,W=S+A \ :\ |A_{jk}|\le t\,S_{jk}\ \bigr\},
\]
et l'on rapporte $\bigl[\inf_{\mathcal{W}}\mathrm{SCR},\ \sup_{\mathcal{W}}\mathrm{SCR}\bigr]$
au lieu d'un point.}

\cheznous{Socle sans contagion $5\,275$~M\euro. Avec contagion et ignorance totale, capital
borné dans $[6\,858\,;\,8\,697]$~M\euro. L'énumération porte sur $1024$ sommets, et un audit
sur $3\,000$ points intérieurs a confirmé qu'aucun ne sort de l'enveloppe des sommets.}

\laphrase{Je ne prétends pas connaître la direction, donc je ne donne pas un chiffre, je
donne une fourchette. Ce n'est pas de la prudence, c'est la seule réponse possible~: une
grandeur qui compte pour le résultat et qui est invisible dans mes données.}

\lepiege{L'atteinte des extrêmes aux sommets est une \emph{propriété vérifiée
numériquement}, pas un théorème démontré. C'est déclaré comme telle dans le mémoire.}

\medskip
\noindent\textbf{Le schéma à projeter si l'on vous demande \og quel est votre SCR\fg{}.}

\begin{center}
\begin{tikzpicture}[xscale=1.0]
\def\y{0}
% axe
\draw[->,draw=navy!60] (-0.4,\y-1.25) -- (11.2,\y-1.25) node[right,font=\scriptsize,text=navy]{M\euro};
\foreach \x/\v in {0/5275, 3.6/6858, 7.4/8697} {
  \draw[draw=navy!60] (\x,\y-1.15) -- (\x,\y-1.35);
  \node[font=\scriptsize,text=navy] at (\x,\y-1.62) {$\v$};
}
% socle
\fill[vert!25,draw=vert!70!black] (0,\y-0.35) rectangle (3.6,\y+0.35);
\node[font=\scriptsize,align=center,text=vert!35!black] at (1.8,\y) {\textbf{socle}\\sans contagion};
% bande d'ignorance
\fill[accent!18,draw=accent!80!black] (3.6,\y-0.35) rectangle (7.4,\y+0.35);
\node[font=\scriptsize,align=center,text=accent!70!black] at (5.5,\y)
  {\textbf{bande d'ignorance}\\largeur $1\,839$};
% resserrement
\draw[decorate,decoration={brace,amplitude=4pt},draw=navy] (7.4,\y+0.5) -- (5.1,\y+0.5)
  node[midway,above=4pt,font=\scriptsize,text=navy,align=center]{un registre spécifié\\en retire $66\,\%$};
\draw[densely dashed,draw=navy!70] (5.1,\y+0.36) -- (5.1,\y+0.5);
\node[align=left,font=\scriptsize,text=navy,anchor=west] at (7.9,\y+0.15)
  {$1\,024$ sommets\\énumérés};
\end{tikzpicture}
\end{center}

\noindent La largeur de la bande n'est pas une incertitude subie~: c'est une grandeur
\textbf{calculée à l'avance}, et le mémoire dit combien elle se réduirait si telle donnée
existait. C'est ce qui transforme une limite en livrable.

\section{Pourquoi $A$ compte quand même}
\label{sec:trace}

\lequoi{Si $A$ est invisible dans la donnée, pourquoi s'en soucier ?}

\laformule{Parce que le capital passe par $(I-W)^{-1}=I+W+W^2+\cdots$, et dès le terme
quadratique $W^2=S^2+SA+AS+A^2$ fait apparaître $A$. En prenant la trace~:
\[
  \operatorname{tr}(W^2)\;=\;\|S\|^2-\|A\|^2 .
\]
Le signe moins vient de l'antisymétrie~: $A^{\top}=-A$ implique $A^2=-AA^{\top}$.}

\laphrase{La direction est invisible au premier ordre mais elle apparaît dès qu'on met la
matrice au carré, et le calcul du capital passe par là. Une grandeur qui compte pour la
réponse et qui est absente de la donnée~: c'est la définition exacte d'une identification
partielle.}

% =====================================================================
\part*{VII. À quel point on y croit}
\addcontentsline{toc}{part}{VII. À quel point on y croit}

\section{Le bootstrap}

\lidee{On ne connaît pas la loi de l'estimateur. Alors on la fabrique~: on rééchantillonne
les données avec remise, on réestime à chaque fois, et la dispersion des estimations
obtenues approche la dispersion vraie.}

\cheznous{Le seul intervalle de confiance bootstrap sur le quantile à $99{,}5\,\%$ vaut un
\textbf{facteur $2{,}6$} entre ses deux bouts. Le bootstrap est aussi utilisé en version
\emph{paramétrique} pour les tests d'adéquation~: on simule sous la loi ajustée pour obtenir
la loi de la statistique de test.}

\laphrase{Le bootstrap fabrique la loi de mon estimateur en rééchantillonnant mes propres
données. Et le résultat est brutal~: rien que cette incertitude-là vaut un facteur deux et
demi sur le capital.}

\lepiege{Le bootstrap capture l'incertitude \emph{de paramètre} à modèle fixé. Il ne dit
rien sur le risque que le modèle lui-même soit faux, qui est bien plus grand ici (facteur
$5{,}5$).}

\section{Les tests d'adéquation}

\lidee{Une loi ajustée peut être fausse. On teste la distance entre la loi ajustée et les
données, et on regarde si cette distance est anormalement grande.}

\laformule{Kolmogorov-Smirnov mesure l'écart maximal entre les fonctions de répartition.
Anderson-Darling pondère ce même écart pour donner davantage de poids aux \emph{queues}, ce
qui est exactement la région qui nous intéresse.}

\cheznous{$p=0{,}87$ pour les deux, par bootstrap paramétrique. La couverture des
intervalles de confiance à $90\,\%$ est de $86\,\%$, donc légèrement optimiste mais
acceptable. La négative binomiale bat le Poisson avec $p\approx10^{-30}$.}

\laphrase{Les tests ne rejettent pas, avec une valeur $p$ de $0{,}87$. Attention à ce que
cela veut dire~: les \emph{formes} sont validées, mais c'est le \emph{niveau} qui reste une
bande.}

\lepiege{Les valeurs critiques tabulées de KS et AD sont fausses quand les paramètres sont
\emph{estimés} sur les mêmes données. Il faut passer par un bootstrap paramétrique, ce que
nous faisons, et il faut le dire, parce que beaucoup de mémoires ne le font pas.}

\section{Les trois étages d'incertitude}

\lidee{L'incertitude ne se réduit pas à l'intervalle de confiance. Elle a trois sources qui
se cumulent et dont les tailles n'ont rien à voir.}

\cheznous{
\begin{center}\small
\begin{tabular}{@{}l l r@{}}
\toprule
\textbf{Étage} & \textbf{Nature} & \textbf{Facteur} \\
\midrule
Paramètre & les données sont finies (bootstrap) & $2{,}6$ \\
Identification & la direction n'est pas mesurable (bornes) & $27\,\%$ de bande \\
Modèle & la famille de loi elle-même est un choix & $5{,}5$ \\
\midrule
\multicolumn{2}{@{}l}{dont la structure de dépendance seule} & $1{,}8$ \\
\bottomrule
\end{tabular}
\end{center}}

\laphrase{Le choix de la famille de queue déplace le résultat d'un facteur cinq et demi,
soit deux fois plus que l'incertitude statistique. Autrement dit, l'essentiel du risque
n'est pas dans mes données, il est dans mes hypothèses. Mais le résultat qui tient est
ailleurs~: \textbf{le niveau bouge dans la bande, l'ordre des piliers ne bouge pas}.}

\lepiege{Cette section peut donner l'impression que rien ne tient. Toujours l'enchaîner
immédiatement avec les tests d'adéquation et avec l'invariance de l'ordre, sinon
l'interlocuteur retient \og flou\fg{} au lieu de \og rigueur\fg{}.}

\section{VaR prédictive et VaR robuste}

\lequoi{La remarque de Caroline Hillairet~: prendre un quantile d'un objet incertain est
mauvais. Comment y répondre sans abandonner la VaR, qui est imposée ?}

\lidee{Deux réponses différentes, à ne pas confondre.
\textbf{La VaR prédictive} intègre l'incertitude de paramètre dans la loi elle-même~: on
mélange les lois sur la loi a posteriori des paramètres, puis on prend le quantile du
mélange.
\textbf{La VaR robuste} prend le pire cas sur un ensemble d'ambiguïté~: on ne mélange pas,
on borne.}

\laformule{
\[
  \mathrm{VaR}^{\text{préd}}_\alpha = \mathrm{VaR}_\alpha\!
  \left(\int f(\cdot\mid\theta)\,\pi(\theta)\,d\theta\right),
  \qquad
  \mathrm{VaR}^{\text{rob}}_\alpha = \sup_{\theta\in\Theta}\ \mathrm{VaR}_\alpha
  \bigl(f(\cdot\mid\theta)\bigr).
\]}

\cheznous{Chaîne complète~: point $659$, prédictive $646$, robuste à $95\,\%$ $1\,026$,
pire cas de modèle $1\,311$. \textbf{La prédictive coïncide presque avec le plug-in} à
$99{,}5\,\%$ ($657$ contre $659$), ce qui est un résultat en soi~: à ce niveau de quantile,
l'intégration ne change presque rien, la fragilité est ailleurs.}

\laphrase{Caroline a raison, et c'est mesurable. Ma réponse n'est pas d'abandonner la VaR,
qui est imposée, mais de rendre l'incertitude visible sur trois étages. Et j'ai un résultat
un peu contre-intuitif~: la VaR prédictive coïncide presque avec la VaR ponctuelle. Ce n'est
pas l'intégration qui manquait, c'est la largeur de bande qui est le vrai sujet.}

\lepiege{Ne pas présenter la prédictive comme \og la bonne réponse\fg{} qui remplacerait le
plug-in. Elle règle un problème (l'incertitude de paramètre) et pas l'autre (l'ambiguïté de
modèle). C'est la robuste qui règle le second, et les deux ne se substituent pas.}

% =====================================================================
\part*{VIII. Les outils annexes}
\addcontentsline{toc}{part}{VIII. Les outils annexes}

\section{Le kappa de Cohen}

\lequoi{Deux codeurs lisent les mêmes récits. Comment mesurer leur accord, sachant que même
deux personnes répondant au hasard seraient parfois d'accord par chance ?}

\laformule{
\[
  \kappa = \frac{p_o - p_e}{1 - p_e},
\]
où $p_o$ est l'accord observé et $p_e$ l'accord attendu par pur hasard. $\kappa=1$ est un
accord parfait, $\kappa=0$ un accord au niveau du hasard.}

\cheznous{Seuil pré-enregistré~: $\kappa\ge0{,}60$ lève la réserve sur le codage. En dessous,
c'est un résultat publiable en soi~: cela signifierait que la séquence des défaillances
n'est pas lisible de façon reproductible dans les rapports d'enquête.}

\laphrase{Le kappa corrige l'accord du hasard. Deux personnes qui répondraient au hasard
seraient d'accord une fois sur cinq~: le kappa retranche cette part. Et j'ai fixé mon seuil
\emph{avant} de connaître le résultat, pour ne pas pouvoir l'ajuster après.}

\section{La valeur d'information}

\lidee{Une recommandation de reporting ne vaut rien tant qu'elle n'est pas chiffrée. On la
chiffre en mesurant de combien la bande de capital rétrécit quand on acquiert
l'information.}

\laformule{Avec $L(\mathcal{V})=\sup_{\mathcal{V}}\mathrm{SCR}-\inf_{\mathcal{V}}\mathrm{SCR}$,
\[
  V(\mathcal{I}) = 1-\frac{L(\mathcal{W}_{\mathcal{I}})}{L(\mathcal{W})} \in[0,1].
\]
Cette quantité est \emph{monotone} (une information plus riche ne vaut pas moins) et
\emph{sans dimension} (donc comparable d'une source de sévérité à l'autre).}

\cheznous{La lecture des sept rapports lève $66\,\%$ de l'ambiguïté. Et la valeur se
\textbf{hiérarchise}~: documenter la seule dépendance gouvernance vers incidents lèverait
$28{,}9\,\%$ à elle seule, tandis que le lien tests vers prestataires n'en lèverait
\emph{rien}.}

\laphrase{J'ai chiffré ce que vaut, en capital, une donnée qu'on n'a pas. Sept rapports
publics lèvent deux tiers de l'ambiguïté. Et surtout un superviseur qui ne pourrait exiger
qu'une seule chose sait laquelle exiger, parce que l'écart entre les dépendances est
énorme.}

\lepiege{Une bande subsisterait même les dix sens connus, imputable à la seule
\emph{amplitude}. L'identification partielle n'est pas un état provisoire que la donnée
finirait par lever, elle est structurelle.}

\section{Le bayésien conjugué}

\lidee{Plutôt que d'estimer un signe de direction par oui ou non, on lui attribue une
probabilité, qu'on met à jour avec les observations documentaires.}

\laformule{Modèle Beta-binomial~: si $\theta\sim\mathrm{Beta}(\alpha_0,\beta_0)$ est la
probabilité que la direction aille dans un sens, et qu'on observe $n_{lo}$ transitions dans
ce sens et $n_{hi}$ dans l'autre, alors
\[
  \theta \mid \text{données} \ \sim\ \mathrm{Beta}(\alpha_0+n_{lo},\ \beta_0+n_{hi}),
\]
et l'on pose $A = S\,(2\theta-1)$, de sorte que $\theta=1/2$ donne $A=0$.}

\cheznous{Médiane $8\,127$~M\euro, intervalle crédible $[7\,957\,;\,8\,316]$. La sensibilité
au prior est de $0{,}8\,\%$, donc négligeable~: les données dominent. Mais seulement
$2$~paires sur $6$ ont un intervalle qui exclut $0{,}50$, donc quatre restent indécises.}

\laphrase{Le bayésien me permet de dire \og ce lien va probablement dans ce sens\fg{} au lieu de
trancher par oui ou non. Deux résultats~: le prior n'a presque aucune influence, ce qui est
rassurant, et seulement deux paires sur six sont réellement tranchées, ce qui est honnête.}

\lepiege{Le bayésien ne crée pas d'information. Il exprime autrement la même incertitude, de
façon plus lisible et plus fine, mais il ne remplace pas la donnée manquante.}

% =====================================================================
\part*{IX. Du capital à la décision}
\addcontentsline{toc}{part}{IX. Du capital à la décision}

\section{À quelle échelle se lit un SCR}

\lequoi{Un capital de plusieurs milliards paraît absurde pour une compagnie d'assurance. C'est
la première objection qu'on vous fera, et une clause de style n'y répond pas.}

\lidee{Le niveau ne dépend pas de la cascade, il dépend de la \textbf{fréquence d'amorce}. On le
teste : on redescend $\lambda$ au niveau d'une entité, à sévérité, matrice $W$ et moteur
\textbf{inchangés}, et l'on regarde ce qui sort.}

\cheznous{\;\\[-6pt]
\begin{center}\small
\begin{tabular}{@{}l r r r@{}}
\toprule
\textbf{Échelle} & $\lambda$/an & \textbf{SCR} & \textbf{Moyenne} \\
\midrule
Secteur financier mondial & $21{,}56$ & $8\,123$~M\euro & $982$~M\euro \\
Seuil \og{}$\geq 10$ ev.\fg{} & $0{,}210$ & $488$~M\euro & $9{,}7$~M\euro \\
Entité moyenne & $0{,}10$ & $221$~M\euro & $4{,}5$~M\euro \\
\bottomrule
\end{tabular}
\end{center}}

\laphrase{Mes montants sont un SCR de \emph{marché}, pas d'entreprise. Redescendu à l'échelle
d'une entité, le même modèle donne $169$ millions, à comparer à la charge opérationnelle de
Formule Standard que je calcule par ailleurs à $450$. Attention au piège que j'ai moi-même
tendu : en lisant $\lambda$ sur les firmes à au moins dix événements, je trouvais $488$ et je
croyais l'écart refermé. Or ces firmes sont dix-neuf fois plus grosses que mon entité. La
lecture correcte donne $0{,}38$ fois le forfait, et ce n'est pas le niveau qui parle, c'est
que le forfait ne bouge pas d'un euro entre un profil exemplaire et un profil défaillant.}

\lepiege{La non-linéarité est le vrai résultat. Diviser $\lambda$ par $235$ divise la moyenne
par $270$, donc à proportion, mais ne divise le capital que par $48$. C'est le
\textbf{principe du grand saut unique} : à $\lambda=0{,}21$, quatre années sur cinq sont vides
et le quantile ne décrit plus une accumulation mais un sinistre isolé. Diviser le capital de
secteur par le nombre d'entités le sous-estimerait d'un facteur six.}

\section{Détenir ou transférer}

\lequoi{Le mémoire compare son SCR à la Formule Standard, aux pertes réelles et au marché
cyber. Trois comparaisons de \emph{niveau}. Il en manque une de \emph{prix}.}

\lidee{Un SCR n'est pas une perte, c'est du capital \textbf{immobilisé}, et le capital coûte.
Solvabilité~II chiffre ce coût explicitement. En face, transférer coûte une prime. Les deux
grandeurs sont homogènes, donc comparables, et c'est l'arbitrage que fait un directeur des
risques.}

\laformule{Le coût annuel de détention vaut $\mathrm{CoC}\times\mathrm{SCR}$, où
$\mathrm{CoC}$ est le taux de coût du capital de la marge de risque : $6\,\%$ historiquement,
ramené à $\mathbf{4{,}75\,\%}$ par la directive (UE)~2025/2. Pour un traité en excès de perte
annuel de portée $L$, on minimise
\[
  \mathrm{co\hat ut}(L)=\underbrace{\E[\min(X,L)]/0{,}485}_{\text{prime}}
  \;+\;\mathrm{CoC}\times\underbrace{\mathrm{VaR}_{99{,}5\%}\bigl(X-\min(X,L)\bigr)}_{\text{capital résiduel}} .
\]}

\cheznous{Multiple de capital $\mathrm{SCR}/\E[X]=\mathbf{47}$ à l'échelle entité contre $8{,}3$
au secteur. Coût de détention $8{,}0$~M\euro/an contre $3{,}6$~M\euro{} de sinistralité
attendue, soit $2{,}2$ fois. Le capital résiduel vaut $\max(\mathrm{SCR}-L,0)$, donc l'optimum
tombe exactement en $L^\star=\mathrm{SCR}=169$~M\euro, et le coût total passe à
$\mathbf{2{,}9}$~M\euro, soit $\mathbf{-64\,\%}$.}

\laphrase{Immobiliser le capital coûte deux fois et demie ce que le risque coûte en moyenne.
Près de l'optimum, un euro de portée coûte un peu plus d'un pour cent de prime et libère
quatre virgule sept cinq pour cent de coût du capital. C'est pour ça que le transfert domine, et
c'est un calcul, pas une opinion.}

\lepiege{Ce gain est une \textbf{borne supérieure}, pour trois raisons : la capacité du marché
plafonne près de $500$~M\euro{} sur un seul risque ; les polices excluent la guerre et les
infrastructures essentielles ; et sous Solvabilité~II céder ne supprime pas l'exigence, cela
\emph{substitue} au risque cyber un risque de \textbf{défaut de contrepartie}. En résumant les
trois par un coefficient d'inefficacité $\kappa$, le seuil critique vaut
$\kappa^\star=1-0{,}0116/0{,}0475=\mathbf{76\,\%}$ : il faudrait que plus des trois quarts de la
couverture soient inefficaces pour renverser la conclusion. La conclusion est robuste, le
chiffre ne l'est pas.}

% =====================================================================
\newpage
\part*{X. Les questions du jury, et quoi répondre}
\addcontentsline{toc}{part}{X. Les questions du jury}

\begin{cle}
\textbf{Comment lire cette partie.} Chaque encadré porte une question que le jury peut poser,
suivie de la réponse à tenir et de la page du mémoire sur laquelle s'appuyer. Les pages
renvoient au PDF déposé, $204$~pages.

\textbf{Trois règles pour l'oral}, tirées du § 6.7 des Recommandations Jury~: la présentation
dure \textbf{25 minutes maximum} et le dépassement est pénalisant~; le jury ne cherche pas à
déstabiliser mais à mesurer la maturité~; et \emph{si vous ne savez pas, dites-le}. Le
référentiel l'écrit noir sur blanc~: mieux vaut dire qu'on ne sait pas que faire semblant.
\end{cle}

\section{Le chiffre et sa nature}

\begin{question}{\og Alors, quel est votre SCR ?\fg{}}
\rep{Le niveau n'est pas une mesure, et le mémoire le démontre plutôt que de s'en excuser~: il
bouge d'un facteur $41$ selon la seule source de données. Ce que le modèle détermine, c'est un
intervalle~: entre $6\,858$ et $8\,697$~M\euro{} sur les $1\,024$ sommets de l'ensemble
admissible. La largeur, $1\,839$~M\euro, est le coût de l'ignorance de la direction, et il est
connu d'avance.}
\appui{résumé p.~5, chapitre~9 p.~108--113.}
\end{question}

\begin{question}{\og Pourquoi refusez-vous d'appeler cela un SCR ?\fg{}}
\rep{Parce qu'il n'existe aucun module DORA en Formule Standard. La grandeur calculée est un
besoin de capital au sens de l'ORSA, donc de pilier~2. L'appeler SCR réglementaire serait
revendiquer une place qu'elle n'a pas dans le dispositif.}
\appui{chapitre~10 p.~114--133, chapitre~2 p.~38--40.}
\end{question}

\begin{question}{\og Vous dites ne pas connaître la direction, et vous publiez un facteur $3{,}34$. Laquelle est votre réponse ?\fg{}}
\rep{Les deux, et elles ne mesurent pas la même chose. Le facteur $3{,}34$ mesure l'effet des
quatre canaux \emph{à direction fixée}~; la bande mesure ce que l'ignorance de la direction
coûte. Le chapitre~10 réconcilie explicitement son SCR conforme de $5\,900$~M\euro{} et le socle
de $5\,275$~M\euro{} du chapitre~9~: ce ne sont pas deux mesures du même objet.}
\appui{chapitre~9 p.~108--113, chapitre~10 p.~114.}
\end{question}

\begin{question}{\og Ce montant sectoriel, que vaut-il pour mon entité ?\fg{}}
\rep{Le mémoire déclare lui-même l'entité notionnelle \emph{hors domaine de validité}~: la
charge décroît beaucoup moins vite que la taille, un facteur $134$ sur les actifs ne divise la
sévérité que par $1{,}53$. Le $169$~M\euro{} est illustratif. Ce qui transfère à une entité,
c'est le rapport au socle et le classement des piliers, pas le niveau.}
\appui{chapitre~10 p.~114--133.}
\end{question}

\begin{question}{\og Pourquoi une VaR et pas une TVaR ?\fg{}}
\rep{Parce que Solvabilité~II définit le SCR comme une VaR à $99{,}5\,\%$ à horizon un an, et
qu'un travail qui se compare au dispositif prudentiel doit parler sa mesure. La TVaR serait plus
cohérente au sens de la théorie du risque, et le mémoire ne prétend pas le contraire~: c'est un
choix de comparabilité, assumé.}
\appui{partie~IV de ce document, chapitre~2 p.~38--40.}
\end{question}

\section{Le modèle et ses hypothèses}

\begin{question}{\og Qu'est-ce qui vous dit que la conformité réduit la propagation ?\fg{}}
\rep{Rien dans les données, et c'est écrit. C'est la seule hypothèse qualitative qui reste, et
le règlement lui-même prend parti puisqu'il impose ces exigences. Mais la thèse n'en dépend pas~:
le capital est croissant en ce paramètre, propriété \emph{démontrée} (monotonie stochastique),
de sorte que toute correspondance qui respecte l'\emph{ordre} reproduit l'écart. La sensibilité
est bornée par construction, pas estimée.}
\appui{chapitre~7 p.~88--94, chapitre~6 p.~79--87.}
\end{question}

\begin{question}{\og Pourquoi une cascade et pas une copule ?\fg{}}
\rep{La copule a été testée en benchmark adverse, pas écartée par principe. Elle sait imiter la
photo d'un état, mais ni le surcoût dû à la propagation, ni l'interaction entre canaux, ni le
durcissement de la dépendance quand l'entité devient non conforme. La dépendance n'est pas le
mécanisme, c'en est l'ombre.}
\appui{chapitre~10 p.~114--133.}
\end{question}

\begin{question}{\og Votre cascade peut-elle exploser ?\fg{}}
\rep{Non, et c'est garanti plutôt qu'espéré. La normalisation de Leontief impose
$\rho(W)\le g<1$, donc la série $I+W+W^2+\cdots$ converge. Le piège est documenté~: la matrice
d'expert brute a un rayon spectral de $1{,}461$, et la brancher telle quelle aurait produit
$21$~entrées négatives dans $(I-W)^{-1}$. C'est pour cela qu'on renormalise.}
\appui{chapitre~6 p.~79--87.}
\end{question}

\begin{question}{\og Vous confondez deux rayons spectraux.\fg{}}
\rep{Non, ils sont distingués explicitement. $\rho(W)$ gouverne le système latent simultané et
doit être $<1$ pour que $(I-W)^{-1}$ existe. $R_0=\rho(M)$ gouverne la cascade d'incidents et
s'éteint si et seulement si $R_0<1$~: c'est le $R_0$ des épidémiologistes. Ce sont deux objets
et deux conditions.}
\appui{chapitre~6 p.~79--87.}
\end{question}

\begin{question}{\og Pourquoi un Vasicek, qui vient du risque de crédit ?\fg{}}
\rep{Parce que la structure est la même~: une variable latente franchit un seuil, et un facteur
commun crée la dépendance. Le mémoire est explicite sur ce qu'il emprunte~: le Vasicek
multifacteur à matrice de corrélation générale est celui de Li (2000), employé tel quel, et il
le déclare comme \emph{n'étant pas} une contribution. La contribution commence à la contagion
dirigée.}
\appui{chapitre~6 p.~79--87.}
\end{question}

\begin{question}{\og Pourquoi avoir abandonné la non-transitivité, qui donnait son titre au mémoire ?\fg{}}
\rep{Parce qu'elle a été testée dans trois formulations et qu'aucune ne tient. La réfuter était
plus utile que la maintenir, et le mémoire le dit au lieu de faire silence. C'est une réponse à
donner franchement~: un travail qui réfute sa propre hypothèse de départ est plus crédible, pas
moins.}
\appui{chapitre~6 p.~79--87.}
\end{question}

\section{Les données}

\begin{question}{\og Votre fréquence vient d'une base américaine sans seuil de matérialité.\fg{}}
\rep{Exact, et c'est pour cela que $\lambda_{\mathrm{ref}}$ est une \textbf{clé de répartition},
pas le niveau du modèle. Le niveau vient d'ailleurs. Le mémoire consacre une section entière à
distinguer les trois grandeurs qui portent le nom de \og fréquence\fg{}, précisément pour éviter
cette confusion.}
\appui{chapitre~4 p.~47--67, chapitre~5 p.~69--78.}
\end{question}

\begin{question}{\og Une de vos sources n'est pas reproductible.\fg{}}
\rep{Hackmageddon, oui, et c'est déclaré comme citation et non comme sortie du dispositif. C'est
aussi la source dont le modèle dépend le moins~: trois instruments indépendants montrent que la
catégorie d'un incident ne porte pas de segmentation ($V$ de Cramér sous $0{,}20$,
$\varepsilon^2$ sous $0{,}08$, rangs sous $0{,}16$). Elle fixe une structure, jamais un niveau.}
\appui{chapitre~4 p.~47--67.}
\end{question}

\begin{question}{\og Comment avez-vous traité les biais de vos bases ?\fg{}}
\rep{En les mesurant plutôt qu'en les déclarant. Deux corrections ont été appliquées à OpRisk et
elles vont en sens opposés~: $0{,}861\times1{,}147=0{,}987$, soit un effet net de $1{,}3\,\%$ sur
la perte moyenne, inférieur à l'incertitude d'estimation de chacune. Ce résultat n'était pas
prévisible et il est publié tel quel.}
\appui{chapitre~4 p.~47--67.}
\end{question}

\begin{question}{\og $582$ pertes, c'est peu pour calibrer une queue.\fg{}}
\rep{C'est peu, et c'est pourquoi la queue est l'incertitude dominante du travail, nommée comme
telle. Sur les $582$, $91$ dépassent le seuil de $20{,}03$~M\euro{} et servent l'ajustement GPD.
Le mémoire ne cache pas cette étroitesse~: il la chiffre en intervalle de confiance sur $\xi$ et
la propage jusqu'au capital.}
\appui{chapitre~5 p.~69--78, chapitre~11 p.~135--146.}
\end{question}

\begin{question}{\og Y a-t-il un risque de double comptage entre vos briques ?\fg{}}
\rep{C'est la question qu'il faut se poser, et le mémoire y répond par la définition des quatre
canaux~: chacun agit sur un objet distinct du modèle (le nombre d'incidents, la détection, la
propagation, l'accumulation par un prestataire commun). L'interaction entre canaux est mesurée
séparément, $5\,001$~M\euro{} soit $35\,\%$ de l'écart, précisément pour ne pas la confondre avec
la somme des effets isolés.}
\appui{chapitre~7 p.~88--94, chapitre~10 p.~114--133.}
\end{question}

\section{L'identification, le c\oe ur du mémoire}

\begin{question}{\og Pourquoi ne pas simplement demander à un expert ?\fg{}}
\rep{C'est ce qui était prévu, et le protocole de Cooke est rédigé et étalonné à l'annexe~C. Il
n'a pas été exécuté faute de panel, et le mémoire le dit dans son titre même~: \og préparé et non
exécuté\fg{}. La lecture centrale repose donc sur un classeur qualitatif, pas sur une agrégation
pondérée d'experts. C'est une limite déclarée, et la première piste léguée.}
\appui{annexe~C p.~179--188, annexe~F p.~194--196.}
\end{question}

\begin{question}{\og Votre corpus de rapports sur-attribue à la gouvernance, et vous concluez qu'elle domine.\fg{}}
\rep{C'est l'objection la plus juste, et elle est dans le mémoire avant d'être dans votre bouche~:
une enquête officielle est \emph{mandatée} pour trouver une cause racine. Mais le biais est
mesuré, pas seulement déclaré. En dégradant l'émission de P1, le classement tient jusqu'au quart
et ne bascule qu'à la moitié~: il faudrait supposer la moitié des arêtes artefactuelles. Le
niveau, lui, ne perd que $24{,}8\,\%$ même à dégradation totale, soit $0{,}81$ fois la largeur de
bande déjà publiée.}
\appui{chapitre~8 p.~95--106, chapitre~10 p.~114--133.}
\end{question}

\begin{question}{\og Un seul codeur a lu vos rapports.\fg{}}
\rep{Oui, et le second codage en aveugle est outillé et attend sa donnée~: le script d'accord
inter-codeurs ne fabrique rien sans le second jeu. C'est la troisième piste de l'annexe~F. En
attendant, le biais de narration est mesuré et son coût en capital recalculé, ce qui vaut mieux
qu'une déclaration d'intention.}
\appui{chapitre~8 p.~95--106, annexe~F p.~194--196.}
\end{question}

\begin{question}{\og Votre expérience naturelle MOVEit, qu'a-t-elle donné ?\fg{}}
\rep{Rien, et c'est un résultat. La différence de différences brute était grande et
significative, $+1{,}17$ jour-brèche avec un intervalle qui excluait zéro. Mais un placebo calé
sur une fausse date un an plus tôt produisait le même effet. La conclusion a donc été retirée.
C'est le genre de test qu'on ne publie que si l'on est prêt à en accepter le verdict.}
\appui{chapitre~8 p.~95--106.}
\end{question}

\begin{question}{\og Votre bande est-elle informative, ou si large qu'elle n'apprend rien ?\fg{}}
\rep{Elle est informative, et la comparaison le montre~: à l'échelle d'entité, la bande
d'identification mesure $51{,}7$~M\euro, quand faire varier le seul indice de queue sur son
intervalle de confiance à $90\,\%$ déplace le capital d'autant. Ignorer la direction coûte donc
autant qu'ignorer la queue, ce qui situe l'enjeu.}
\appui{chapitre~10 p.~114--133.}
\end{question}

\section{La validation}

\begin{question}{\og On ne peut pas valider un quantile à $99{,}5\,\%$.\fg{}}
\rep{Exact, et le mémoire le chiffre au lieu de l'éluder~: détecter un taux de dépassement
\emph{double} du nominal avec une puissance de $80\,\%$ demanderait $1\,811$~années à
$99{,}5\,\%$. D'où le dispositif retenu~: valider la \emph{forme} des lois et la fréquence, et
déclarer le niveau non testable.}
\appui{chapitre~5 p.~69--78.}
\end{question}

\begin{question}{\og Votre backtest rejette le niveau.\fg{}}
\rep{Oui, et pour un motif nommé et mesuré~: une dérive de l'échelle de sévérité de
$4{,}00\,\%$ par an, que la calibration stationnaire n'absorbe pas. La loi de fréquence couvre
$91{,}7\,\%$ des années pour $90\,\%$ annoncés, la forme de queue passe le test de probabilité
intégrale. Le défaut est chiffré plutôt que corrigé, et légué comme piste de mémoire.}
\appui{chapitre~11 p.~135--146, annexe~F p.~194--196.}
\end{question}

\begin{question}{\og Comment savez-vous que votre moteur de simulation ne se trompe pas ?\fg{}}
\rep{Il est contrôlé par trois chemins indépendants, dont une inversion de Fourier de la loi
composée, et l'écart à la simulation reste inférieur à une erreur type. Le mémoire précise aussi
la limite de ce contrôle~: les quatre chemins partagent les lois et les tables de cascade, donc
ils vérifient l'\emph{agrégation}, jamais la vérité du modèle.}
\appui{chapitre~11 p.~135--146.}
\end{question}

\begin{question}{\og $\xi=0{,}595$, donc votre moyenne est infinie ?\fg{}}
\rep{Non. Pour $1/2<\xi<1$, la variance est infinie et l'espérance finie. C'est exactement ce
qu'écrit le mémoire, et la proposition qui donne quantile et moyenne de queue pose explicitement
$\xi\in(0,1)$ dans ses hypothèses.}
\appui{chapitre~5 p.~69--78.}
\end{question}

\section{Le métier et la décision}

\begin{question}{\og Qu'est-ce qu'un dirigeant fait de votre travail lundi matin ?\fg{}}
\rep{Trois choses. Il sait quel pilier remédier en premier, et que l'ordre optimal n'est pas
l'ordre glouton parce que les interactions imposent de raisonner sur l'horizon entier. Il sait
comparer remédier, détenir et transférer sur un seul axe en euros par an. Et il sait quelle
donnée réclamer à son registre d'incidents pour resserrer son propre chiffre.}
\appui{chapitre~10 p.~114--133, chapitre~12 p.~147--149.}
\end{question}

\begin{question}{\og Votre recommandation dépend-elle de paramètres que personne ne sait calibrer ?\fg{}}
\rep{Non, et c'est un résultat à mettre en avant~: l'ordre de remédiation optimal ne dépend que
des SCR par configuration, donc il est \emph{invariant} aux taux de transition de la chaîne de
Markov, qui sont pourtant les paramètres les moins calibrables du modèle. Ces taux fixent le
calendrier, jamais la séquence.}
\appui{chapitre~10 p.~114--133.}
\end{question}

\begin{question}{\og Que répondez-vous à un superviseur ?\fg{}}
\rep{Une recommandation chiffrée et opposable~: rendre obligatoires deux champs \emph{déjà
normalisés} dans le schéma VERIS, via la notification d'incident majeur sous DORA. Le coût de
conception est nul puisque la taxonomie existe. Le gain est mesuré~: $66\,\%$ de resserrement de
bande, et $28{,}9\,\%$ pour la seule dépendance gouvernance-incidents.}
\appui{chapitre~12 p.~147--149.}
\end{question}

\begin{question}{\og Votre travail apporte quoi à la science actuarielle ?\fg{}}
\rep{Un mécanisme de dépendance \emph{dirigée} dont trois propriétés sont démontrées, dont
celle-ci~: une matrice et sa transposée, indiscernables pour la donnée, donnent un capital et une
décision différents. Une frontière d'identification établie sur données réelles. Et une méthode
importée de l'économétrie, l'identification partielle, appliquée à une exigence de capital.}
\appui{conclusion p.~150--153.}
\end{question}

\section{La forme, la déontologie, le dispositif}

\begin{question}{\og Avez-vous utilisé l'intelligence artificielle générative ?\fg{}}
\rep{Oui, et c'est déclaré en fin de résumé et détaillé à l'annexe~G, comme le référentiel
l'exige. Assistance à la rédaction, écriture et relecture des scripts, agents de vérification.
Aucun résultat chiffré n'en sort~: tout nombre publié vient de l'exécution d'un script du dépôt,
et un dispositif rattache chaque nombre à la sortie qui le produit. L'annexe dit aussi ce que ce
dispositif ne détecte pas.}
\appui{résumé p.~5, annexe~G p.~197--199.}
\end{question}

\begin{question}{\og Comment garantissez-vous qu'aucun chiffre n'est faux ?\fg{}}
\rep{Je ne le garantis pas, et l'annexe~G explique pourquoi. Le harnais confronte $1\,839$
nombres publiés aux sorties des scripts cités, et il est à $100\,\%$. Mais une relecture au filtre
des superlatifs et des comptages a trouvé huit erreurs sémantiques que le harnais confirmait
toutes les huit~: le nombre sortait du bon script, seule la phrase autour était fausse. Un taux de
confirmation ne mesure pas la justesse du texte.}
\appui{annexe~G p.~197--199, conclusion p.~150--153.}
\end{question}

\begin{question}{\og Votre mémoire fait $204$ pages, le référentiel en recommande $70$ de corps.\fg{}}
\rep{Le référentiel écrit \og environ\fg{} et ne prévoit aucune sanction. Le volume vient des données
et de la robustesse~: $21$~pages de données, $12$~de robustesse, et $100\,\%$ des nombres publiés
rattachés à un script. Rien n'y est du remplissage, et le matériel qui pouvait sortir est sorti,
vers un dépôt Zenodo pérenne.}
\appui{notice des annexes p.~160.}
\end{question}

\begin{question}{\og Pourquoi avoir déporté des annexes hors du document ?\fg{}}
\rep{Pour alléger la lecture sans perdre la traçabilité. Les démonstrations propres à la
contribution, celles de la cascade, restent dans le document~; ce qui est sorti, ce sont les
théorèmes standards des valeurs extrêmes, les pièces justificatives et les compléments. Le tout
est déposé sous DOI, donc citable et pérenne.}
\appui{notice p.~160, annexe~A p.~161--165.}
\end{question}

\begin{question}{\og Qu'auriez-vous fait avec six mois de plus ?\fg{}}
\rep{Exécuter l'élicitation, qui est prête. Faire coder le corpus en aveugle par un second
lecteur. Et attaquer l'amplitude des liens, pas seulement leur sens, parce que le mémoire
démontre qu'une bande subsisterait même les dix directions connues. Les six pistes sont écrites
en annexe~F, avec ce qui existe déjà pour chacune.}
\appui{annexe~F p.~194--196.}
\end{question}

\begin{question}{\og Quelle est la faiblesse de votre travail ?\fg{}}
\rep{À dire sans détour, parce que le jury le sait déjà et que l'esquive coûte plus que
l'aveu~: le niveau absolu du capital n'est pas une mesure, et l'ordre entre les deux piliers de
tête dépend d'une direction que je ne sais pas identifier. Ce que je sais dire, c'est que ces
deux-là dominent les trois autres, et de combien l'ignorance coûte.}
\appui{conclusion p.~150--153.}
\end{question}

\begin{cle}
\textbf{La règle qui vaut pour toutes ces réponses.} Aucune ne consiste à défendre un chiffre~:
toutes consistent à dire ce que le chiffre vaut et ce qu'il ne vaut pas. C'est la posture que le
§ 4.6 du référentiel valorise explicitement, et c'est celle qui a produit ce mémoire.
\end{cle}

% =====================================================================
\newpage
\section*{Annexe~: les nombres à connaître par c\oe ur}
\addcontentsline{toc}{section}{Annexe : les nombres à connaître}

\begin{center}\small
\renewcommand{\arraystretch}{1.25}
\begin{tabular}{@{}l l l@{}}
\toprule
\textbf{Nombre} & \textbf{Ce que c'est} & \textbf{Où} \\
\midrule
\multicolumn{3}{@{}l}{\textcolor{navy}{\textbf{Le chiffre de tête, à savoir sans hésiter}}} \\
$6\,049 \to 20\,188$~M\euro & capital, état conforme vers non conforme & résultat \\
$3{,}34$ & le facteur entre les deux états & résultat \\
$14\,139$~M\euro & l'écart, dont $9\,138$ isolés & résultat \\
$5\,001$~M\euro\ ($35\,\%$) & la part portée par la seule interaction & résultat \\
$[6\,858\,;\,8\,697]$~M\euro & \textbf{la réponse à \og quel est votre SCR\fg{}} & identification \\
$1\,839$~M\euro & largeur de bande = coût de l'ignorance & identification \\
\midrule
$\varphi = 9{,}2$ & surdispersion de la fréquence & fréquence \\
$p\approx10^{-30}$ & rejet de Poisson contre NB & fréquence \\
$\xi=0{,}5954$ / $1{,}03$ & indice de queue OpRisk / PRC & sévérité \\
$u=20{,}03$~M\euro & \textcolor{accent}{seuil des excès (corrigé)} & sévérité \\
$91$ & excès au-dessus du seuil & sévérité \\
$582$ / $27$~ans & pertes OpRisk, périmètre cyber-finance & données \\
$15\,053$ & notifications PRC, $2019$--$2025$ & données \\
$p_u=0{,}150$ & part des pertes au dessus du seuil & sévérité \\
$g=0{,}90$ & facteur de propagation & cascade \\
$\rho(\mathrm{TRANS})=1{,}461$ & pourquoi il faut renormaliser & cascade \\
$\|A\|/\|W\|=40{,}7\,\%$ & part directionnelle de la matrice & cascade \\
$K_4=1{,}033$ & seuil de déclenchement de P4 & cascade \\
$\lambda_U=0{,}53$ & dépendance de queue de Gumbel & dépendance \\
$0{,}17 \to 0{,}39$ & co-extrême, gaussien vers Gumbel & dépendance \\
$119$ vs $128\pm27$ & placebo OpRisk & identification \\
$z=-0{,}33$ & placebo, bases agrégées & identification \\
$z=+3{,}93 \to +5{,}12$ & placebo post-mortems, 7 puis 10 rapports & identification \\
$+4{,}08$ / $+3{,}62$ & le même, primaires seules / entités seules & identification \\
$0{,}12\,\%$ / $1{,}19\,\%$ / $1$ & \textsc{veris} : champ, chronologie, intersection & identification \\
$12$ / $0$ sur $20$ & liens $p_{jk}$ crédiblement faibles / forts & identification \\
$\sigma=0{,}0700$ & production d'entropie & identification \\
$1024$ & sommets énumérés & identification \\
$5\,275$~M\euro & socle sans contagion & capital \\
$[6\,858\,;\,8\,697]$~M\euro & capital borné & capital \\
$8\,123$ / $169$~M\euro & SCR secteur / entité ($\lambda$ lu à la taille) & échelle \\
$50$ contre $8{,}3$ & multiple de capital, entité contre secteur & échelle \\
$0{,}087$ $[0{,}026\,;0{,}148]$ & élasticité sévérité-taille, corrigée du seuil & données \\
$+163\,\%$ & rehaussement chain-ladder de 2023--2025 & données \\
$4{,}75\,\%$ & coût du capital, directive (UE) 2025/2 & décision \\
$-64\,\%$ et $\kappa^\star=78\,\%$ & \textcolor{accent}{gain du traité, seuil de robustesse (corrigé)} & décision \\
$p=0{,}87$ & tests d'adéquation & validation \\
$2{,}6$ / $1{,}8$ / $5{,}5$ & les trois étages d'incertitude & validation \\
$91{,}7\,\%$ / $75{,}0\,\%$ & backtest~: fréquence couverte / niveau & validation \\
$4{,}00\,\%$ par an & la dérive d'échelle qui explique le rejet & validation \\
$1\,811$ années & pour tester $99{,}5\,\%$ à $80\,\%$ de puissance & validation \\
$66\,\%$ / $28{,}9\,\%$ & valeur d'information, totale / meilleure paire & livrable \\
\midrule
\multicolumn{3}{@{}l}{\textcolor{navy}{\textbf{Robustesse du classement, à sortir si l'on attaque le corpus}}} \\
$43{,}6\,\%$ / $37{,}7\,\%$ & P1 / P4 premiers sur l'ensemble admissible & robustesse \\
$1{,}5\,\%$ & écart de regret maximal entre les deux & robustesse \\
$a=0{,}25$ / $a=0{,}50$ & le classement tient / bascule vers P4 & robustesse \\
$-24{,}8\,\%$ ($0{,}81\times$) & chute du niveau si P1 devient muette & robustesse \\
$34\,\%$ & part Shapley de P1 dans le surcoût & résultat \\
\bottomrule
\end{tabular}
\end{center}

\vspace{4mm}
\begin{cle}
\textbf{La phrase qui tient tout.} Si l'on ne devait retenir qu'une chose de ce document~:
\emph{le niveau de capital bouge dans une bande large, mais l'ordre des piliers ne bouge
sous aucune des hypothèses testées}. C'est la thèse du mémoire, et c'est le seul énoncé qui
survit aux trois étages d'incertitude.
\end{cle}

\end{document}
